\subsection{An optional improvement using generic chaining}
The extra sample-size logarithms in Lemma~\ref{lem:unknown-link} can be
removed by using the two-sided majorizing-measure theorem~\cite{talagrand} instead of
Sudakov minoration followed by an entropy integral. We state the
additional imported theorem explicitly: for a compact index set with
its canonical Gaussian metric, the expected supremum of the centered
Gaussian process is bounded above and below, up to absolute constants,
by Talagrand's $\gamma_2$ functional. We also use the generic-chaining
upper bound for a subgaussian process and the entropy-integral upper
bound on $\gamma_2$.

For a symmetric compact $F\subset[-1,1]^n$ containing zero, retain the
metric $d_n$ and normalized Gaussian complexity $G(F)$ above. The scalar
class $\mathcal L$ is compact in the uniform norm by Arzel\`a--Ascoli.
Its entropy estimate gives
\[
 \gamma_2(\mathcal L,\|\cdot\|_\infty)
 \le C\int_0^2\sqrt{\log\mathcal N(a,\mathcal L,\|\cdot\|_\infty)}\,da
 \le C.
\]
Equip $F\times\mathcal L$ with the sum metric
$d_n(f,f')+\|\psi-\psi'\|_\infty$. Admissible chaining nets for the two
factors, shifted by one level before taking their products, give
\[
 \gamma_2(F\times\mathcal L)
 \le C\{\gamma_2(F,d_n)+\gamma_2(\mathcal L,\|\cdot\|_\infty)\}.
\]
The shift is sufficient because the product of two sets of cardinality
at most $2^{2^{k-1}}$ has cardinality at most $2^{2^k}$.
The map $(f,\psi)\mapsto\psi\circ f$ is one-Lipschitz for this sum
metric. The Gaussian process indexed by its image has canonical metric
$d_n/\sqrt n$. Generic chaining and then the majorizing-measure lower
bound for the Gaussian process indexed by $F$ yield
\[
 G(\mathcal L\circ F)
 \le\frac C{\sqrt n}\{\gamma_2(F,d_n)+1\}
 \le C G(F)+C/\sqrt n.
\]
All processes are centered; symmetry and the zero index make the
expected-supremum conventions unambiguous.

Gaussian width is unchanged by the absolute convex hull of a symmetric
set. Adding the constant atom costs at most $1/\sqrt n$.
The recursion for $\mathcal A_H$ therefore gives
\[
 G(\mathcal A_H)
 \le C^{H+1}[1+\sqrt{\log(2d+2)}]/\sqrt n.
\]
Finally $\mathfrak R(F)\le\sqrt{\pi/2}\,G(F)$: condition on Gaussian
signs and apply Jensen to the independent magnitudes, whose common
mean is $\sqrt{2/\pi}$. Thus the factor
$\log^{3H/2}(2n)$ may be omitted from
\eqref{eq:all-lipschitz-full-energy-rad}, at the cost of another absolute
constant to the power $H$. This optional refinement uses a deeper
standard Gaussian-process theorem; the elementary proof above already
establishes all polynomial-learning and activation-scope claims.
